\section{Related Work}
\label{sec:related}

\para{Refusal mechanisms.}
The closest mechanistic analogy to a bedtime nudge is refusal. \citet{arditi2024refusal} show that refusal in several open-source chat models can be mediated by a single activation-space direction, and \citet{wang2025universal} report cross-lingual transfer of a refusal direction across safety-aligned languages. These results motivate the question of whether other assistant boundaries might also be controlled by unified mechanisms. Our study differs in two ways: we evaluate closed-source API behavior rather than internal activations, and the target behavior is not refusal of harmful content but unsolicited rest advice during safe requests.

\para{False refusal and over-refusal.}
Benchmark work shows that models can decline safe prompts when they resemble unsafe ones. \xstest targets exaggerated safety behavior on safe prompts with unsafe-looking surface forms \cite{rottger2023xstest}; \orbench scales over-refusal evaluation and highlights a safety-helpfulness trade-off \cite{cui2024orbench}; \donotanswer and \sorrybench provide broader safety-refusal taxonomies and evaluation settings \cite{wang2023donotanswer,xie2024sorrybench}. Mechanistic work further separates true refusal from false refusal and over-refusal \cite{wang2024surgical,maskey2026overrefusal,joad2026more}. We build on this distinction by treating \bedtimenudges as a separate category: a response can answer the user and still include a sleep suggestion, so simple refusal/compliance labels are not enough.

\para{User characteristics and personalization.}
Persona prompts and user histories can change model behavior. \citet{plazadelarco2025nofor} find that sociodemographic persona prompts can affect false refusal, though model and task effects are often larger. \citet{jain2025interaction} show that interaction context often increases sycophancy, and \citet{kelley2026personalization} find that personalization can increase affective alignment while affecting epistemic independence differently by role. Long-horizon personalization benchmarks such as \realpref and \personamem motivate evaluations that separate user-profile effects from generic context-length effects \cite{guo2026realistic}. We use this line of work to design neutral, fatigued, and dependent profiles while holding the final task safe and concrete.

\para{Vulnerable conversations and real-world chat corpora.}
Companion and vulnerable-conversation studies show that platforms can differ in advice, presence, and corrective friction during extended interactions \cite{chu2026chatbots}. Human-AI interaction studies also show that conversational systems absorb breakdowns and task-specific repair patterns \cite{mahmood2023user}. Real-world corpora such as \wildchat provide a way to audit naturally occurring phrases \cite{wildchat2024}, but phrase matches can conflate user-requested sleep advice, ordinary sleep topics, and assistant-initiated nudges. Our \wildchat audit therefore serves only as observational context; the causal test is the synthetic factorial design in \secref{sec:methodology}.
